% !TeX root = ../main.tex

% ----------------------------------------------------------
% Considerações Finais
% ----------------------------------------------------------
\chapter*[Considerações Finais]{Considerações Finais}
\addcontentsline{toc}{chapter}{Considerações Finais}

Este trabalho avaliou as diferentes abordagens algorítmicas para a resolução de problemas de otimização combinatória nas quais encontrar a solução exata tem custo computacional elevado, como o Problema do Caixeiro-Viajante como caso central e estendendo a análise para outros problemas, como a Cobertura de Vértices, \textit{Set Cover} e \textit{Bin Packing}. A dificuldade de resolver problemas deste tipo se dá pela explosão combinatória: por exemplo, o problema do caixeiro viajante com 20 cidades possui mais de 60,8 quatrilhões de rotas distintas e, além disso, o esforço computacional exponencial ou fatorial não foi superado pela evolução da velocidade dos processadores. Assim, escolher uma estratégia significa, portanto, realizar um \textit{trade-off} entre tempo ou garantia da solução encontrada.

As estratégias exatas retornam sempre a solução ótima do problema, além de provar que a solução encontrada realmente é ótima. A busca exaustiva garante essa propriedade, porém o tempo gasto para isso é alto, e o método \textit{branch-and-bound} reduz este tempo sem perder a garantia: no exemplo de 4 cidades, a rota gulosa de 134~km é utilizada como limite superior inicial, em seguida é substituída pela rota ótima de 99~km, e ramos como $a\text{-}b\text{-}c$ e $a\text{-}b\text{-}d$ são podados, pois os limites inferiores são, respectivamente, 135 e 134~km, sem serem completados. O pior caso ainda assim é exponencial. A diferença entre os métodos \textit{branch-and-bound} e \textit{backtracking} está no critério utilizado para podar os ramos: enquanto um poda ramos viáveis se o melhor valor possível neste ramo não for melhor que a solução incumbente, o outro poda ramos inviáveis. \textbf{(Cap.~3: acrescentar informações)}

Já as heurísticas atuam de forma análoga: abrem mão da prova em favor da velocidade. O vizinho mais próximo constrói uma rota em tempo $O(n^2)$, mas na mesma instância de 4 cidades, ele encontrou uma rota com 134~km, 35\% a mais em relação ao ótimo de 99~km. Isso acontece porque a escolha da melhor opção localmente deixou a aresta de 100~km para o final. Esse comportamento é característico das heurísticas: sem uma prova, olhar apenas o resultado não permite saber se a solução está próxima do ótimo ou muito longe dele. A busca local ameniza esse problema ao melhorar a rota gulosa por pequenas trocas: em cinco cidades, duas aplicações do 2-opt reduziram o percurso de 40,30 para 30,13~km, alcançando o ótimo da instância. Ainda assim, ela só assegura um ótimo local; parar sem melhorias não prova que a rota é globalmente ótima.

Por fim, os algoritmos de aproximação estão localizados entre essas duas famílias: também mantêm o tempo polinomial, mas adicionam à heurística uma prova de que, para qualquer instância, o custo da solução encontrada fica a um fator $\rho$ da solução ótima. Para a Cobertura de Vértices, o algoritmo de emparelhamento maximal é uma 2-aproximação, e a prova dessa afirmação é ilustrada como a ideia central do capítulo~5: como não conhecemos o ótimo, a garantia o substitui por um limite inferior sendo o número de arestas escolhidas. No exemplo do grafo caminho com 6 vértices, nas duas execuções do algoritmo, são retornadas coberturas com tamanho 6 e 4, contra o ótimo de 3, ou seja, razões de aproximação 2 e aproximadamente 1,33 respectivamente, a depender da ordem de escolha, enquanto a razão garantida vale sempre. No problema de \textit{Set Cover}, a razão do algoritmo guloso de custo-benefício é $H(n) \le 1 + \ln n$. Já no exemplo da cobertura máxima orçada, variante com o limite do orçamento, o guloso cobre 9 dos 10 moradores, enquanto a solução ótima cobria todos os 10; isso prova que a escolha gulosa não garante o ótimo. \textbf{(Cap.~6: acrescentar informações).}

Em síntese, as três famílias se distinguem pelo que provam: a exata, a solução ótima, a aproximada, um limite para o erro e a heurística, que não prova nenhum dos casos anteriores. Por isso, a mesma ideia gulosa pode assumir papéis diferentes conforme o problema ao qual é aplicada: enquanto no TSP o vizinho mais próximo é uma heurística sem garantia, no \textit{Set Cover} o guloso de custo-benefício é uma aproximação de razão conhecida.

Portanto, cada metodologia será eficiente a depender do problema em que é aplicada, não existindo uma estratégia universalmente melhor em relação às demais. Problemas pequenos ou aqueles onde a prova da otimalidade é essencial são favorecidos pelas estratégias exatas, já problemas cujas decisões devem ser tomadas em um curto intervalo de tempo, como o despacho diário de uma frota por exemplo, ou em ambientes que mudam o tempo todo são beneficiados por heurísticas e algoritmos de aproximação, já que raramente compensará o custo de refinar a eficiência de uma solução de 98\% para 100\%. Quando se quer saber qual o pior que pode acontecer, a aproximação é a escolha mais natural. Porém, é importante ressaltar que nem todo problema difícil é aproximável de forma eficiente, com sua garantia dependendo das hipóteses do problema como a desigualdade triangular no TSP, dispensável na Cobertura de Vértices. Logo, compreender o que cada estratégia garante e a que custo é tão importante quanto conhecer o algoritmo em si.
% ----------------------------------------------------------
